\scanpage{115}
de sorte que $\zeta_1,\zeta_2,\zeta_3$ sont les solutions de l'équation
\[
\underbrace{z^3+a_1z^2+a_2z+a_3}_{\substack{\deleted{\unsure{(z-\zeta_1)(z-\zeta_2)|(z-\zeta_j)(z-\zeta_k)}}\\\prod_i(z-\zeta_i)}}=0\tag*{(60)}
\]
On aura, \deleted{\unsure{donc|dans}} si
\[
\delta\in t_E=T_{E,0}\quad\deleted{\unsure{\delta|\xi}}
\]
est une base de $t_E$, liée à $p\bmod k.1$ par
\[
\varepsilon_2(p)\simeq\delta^{\otimes2}
\]
et si
\[
p'=\Theta_\delta p
\]
la relation
\[
p'{}^2=p^3+a_1p^{\unsure{2|7}}+a_2p+a_3.
\]

Déterminons l'autom. involutif $\sigma_i$ de $\mathbb{P}^1_k$ qui échange $\zeta_i$ et $\infty$, $\zeta_j$ et $\zeta_k$ -- qui \deleted{\unsure{est donc|est donné}}\inserted{s'identifie} à l'autom. déduit par passage au quotient de
\[
\deleted{x\longmapsto x+}\qquad x\longmapsto x+\varepsilon_i\qquad E\longrightarrow E
\]
On trouve, \deleted{\unsure{utilisant|utilisons}}\inserted{écrivant,}
\[
\sigma_i z=\frac{az\unsure{-|+}b}{cz\unsure{-|+}d}
\]
et écrivant $\sigma_i\infty=\zeta_i,\ \sigma_i\zeta_i=\infty$
\[
\deleted{\sigma_i z}\quad\frac ac=\zeta_i,\quad-\frac dc=\zeta_i,
\]
soit, faisant $c=1$
\[
\sigma_i z=\frac{\zeta_i z+b}{\deleted{\unsure{\zeta_i|a_i}}z-\unsure{\zeta_i|\xi_i}}
\]
et écrivant $\sigma_i\zeta_j=\zeta_k$ on trouve
\[
\frac{\unsure{\zeta_i|\xi_i}\zeta_j+b}{\zeta_j-\zeta_i}=\zeta_k\qquad\text{i.e.}\qquad b=\zeta_j\zeta_k-\zeta_i(\zeta_j+\zeta_k)
\]
% archive sheet number: 114
